\section{Complex series}
	\begin{itemize}
		\item \textit{Power series}: A series in powers of $(z-z_0)$ is known as a power series, i.e.
		\begin{align*}
			\sum_{n=0}^{\infty}{a_n(z-z_0)^n}&=a_0+a_1(z-z_0)+a_2(z-z_0)^2+\cdots+a_k(z-z_0)^k+\cdots
		\end{align*}
		where $z$ is a complex variable and $z_0$ is the centre of the series. The constants $a_0$, $a_1$, $a_2$, $\dots$ are known as the coefficients of the series.
		
		\item \textit{Region of convergence}: The region of convergence for a power series is a set of points over which the series is defined. There may be three distinct possibilities for the region of convergence:
		\begin{enumerate}[i.]
			\item The series converges only at the point $z=z_0$.
			\item The series converges for all points $z$ on the complex plane.
			\item The series converges everywhere inside a circular region given as $|z-z_0|<R$, where $z_0$ is the \textit{centre of convergence} and $R$ is its \textit{radius of convergence}.
		\end{enumerate}
		
		\item \textit{Radius of convergence}: Consider the power series which is centred at the origin, i.e. $\sum_{n=0}^{\infty}{a_nz^n}$. Then its radius of convergence $R$ is given by
		\begin{enumerate}[a)]
			\item $\dfrac{1}{R}=\lim_{n\to\infty}{\qty|a_n|^{1/n}}$
			
			\item $\dfrac{1}{R}=\lim_{n\to\infty}{\qty|\dfrac{a_{n+1}}{a_n}|}$
		\end{enumerate}
	\end{itemize}
	
		\subsubsection*{Examples}
			\begin{enumerate}
				\item Find the radius of convergence of the power series $1+z+\dfrac{z^2}{2!}+\dfrac{z^3}{3!}+\cdots$.\\
				\textbf{Solution:} The given power series can be written as
				\begin{align*}
					\sum_{n=0}^{\infty}{a_nz^n}&=\sum_{n=0}^{\infty}{\dfrac{z^n}{n!}} \\
					\implies a_n&=\dfrac{1}{n!}
				\end{align*}
				Hence the radius of convergence is
				\begin{align*}
					\dfrac{1}{R}&=\lim_{n\to\infty}{\qty|\dfrac{a_{n+1}}{a_n}|} \\
					\implies\dfrac{1}{R}&=\lim_{n\to\infty}{\qty|\dfrac{{1}/{(n+1)!}}{{1}/{n!}}|}=\lim_{n\to\infty}{\qty|\dfrac{n!}{(n+1)!}|}=\lim_{n\to\infty}{\qty|\dfrac{n!}{(n+1)n!}|} \\
					\implies\dfrac{1}{R}&=\lim_{n\to\infty}{\qty|\dfrac{1}{n+1}|}=\dfrac{1}{\infty} \\
					\implies\Aboxed{R&=\infty}
				\end{align*}
				which implies that the given power series converges everywhere on the complex plane.
				
				\item Find the radius of convergence of the power series $\sum_{n=0}^{\infty}{\dfrac{z^n}{2^n+3}}$.\\
				\textbf{Solution:} We have here
				\begin{align*}
					\sum_{n=0}^{\infty}{a_nz^n}&=\sum_{n=0}^{\infty}{\dfrac{z^n}{2^n+3}} \\
					\implies a_n&=\dfrac{1}{2^n+3} \\
					\therefore\;a_{n+1}&=\dfrac{1}{2^{n+1}+3}
				\end{align*}
				
				Hence the radius of convergence is
				\begin{align*}
					\dfrac{1}{R}&=\lim_{n\to\infty}{\qty|\dfrac{a_{n+1}}{a_n}|} \\
					\implies\dfrac{1}{R}&=\lim_{n\to\infty}{\qty|\dfrac{{1}/(2^{n+1}+3)}{{1}/(2^n+3)}|}=\lim_{n\to\infty}{\qty|\dfrac{2^n+3}{2^{n+1}+3}|}=\lim_{n\to\infty}{\qty|\dfrac{2^n(1+3/2^n)}{2^n(2+3/2^n)}|} \\
					\implies\dfrac{1}{R}&=\dfrac{1+3/2^\infty}{2+3/2^\infty}=\dfrac{1}{2} \\
					\implies\Aboxed{R&=2}
				\end{align*}
				which implies that the given power series converges for all points in the circular region $|z|<2$.
				
				\item Find the radius of convergence of the power series $\sum_{n=0}^{\infty}{\dfrac{n!}{n^n}z^n}$.\\
				\textbf{Solution:} We have here
				\begin{align*}
					\sum_{n=0}^{\infty}{a_nz^n}&=\sum_{n=0}^{\infty}{\dfrac{n!}{n^n}z^n} \\
					\implies a_n&=\dfrac{n!}{n^n} \\
					\therefore\;a_{n+1}&=\dfrac{(n+1)!}{(n+1)^{n+1}}
				\end{align*}
				So we obtain
				\begin{align*}
					\dfrac{a_{n+1}}{a_n}&=\dfrac{(n+1)!}{(n+1)^{n+1}}\cdot\dfrac{n^n}{n!} \\
						&=\dfrac{\cancel{(n+1)}\cancel{n!}}{(n+1)^{n}\cancel{(n+1)}}\cdot\dfrac{n^n}{\cancel{n!}}=\dfrac{n^n}{(n+1)^n} \\
						&=\dfrac{\cancel{n^n}}{\cancel{n^n}\qty(1+1/n)^n} \\
					\implies\dfrac{a_{n+1}}{a_n}&=\dfrac{1}{\qty(1+\dfrac{1}{n})^n}
				\end{align*}
				
				Hence the radius of convergence is
				\begin{align*}
					\dfrac{1}{R}&=\lim_{n\to\infty}{\qty|\dfrac{a_{n+1}}{a_n}|} \\
						&=\lim_{n\to\infty}{\qty|\dfrac{1}{\qty(1+\dfrac{1}{n})^n}|} \\
						&=\qty|\dfrac{1}{\lim\limits_{n\to\infty}{\qty(1+\dfrac{1}{n})^n}}|
				\end{align*}
				
				Now, the Euler's constant $e$ can be obtained as a limit as given below
				\begin{align*}
					e=\lim_{n\to\infty}{\qty(1+\dfrac{1}{n})^n}
				\end{align*}
				
				Therefore, we finally obtain the radius of convergence as
				\begin{align*}
					\dfrac{1}{R}&=\qty|\dfrac{1}{e}|=\dfrac{1}{e} \\
					\implies\Aboxed{R&=e}
				\end{align*}
				
				which implies that the given power series converges for all points in the circular region $|z|<e$.
			\end{enumerate}
	
	\subsection{Taylor's series expansion of a complex function}
	
		\subsubsection{Taylor's theorem}
			If a complex function $f(z)$ is analytic at all points inside a circle $C$, with its centre at the point $a$ and radius $R$, then at each point $z$ inside $C$, we can expand $f(z)$ in the form of a Taylor's power series as given below:
			\begin{align*}
				f(z)&=f(a)+f'(a)(z-a)+\dfrac{f''(a)}{2!}(z-a)^2+\dfrac{f'''(a)}{3!}(z-a)^3+\cdots \\
				\implies f(z)&=\sum_{n=0}^{\infty}{\dfrac{f^{(n)}(a)}{n!}(z-a)^n}
			\end{align*}
			where $f^{(n)}(a)=\left.\dv[n]{f}{z}\right|_{z=a}$ is the $n^\text{th}$ derivative of $f(z)$ at $z=a$.
	
	\newpage
	\subsection{Laurent's series expansion of a complex function}
		
		\subsubsection{Laurent's theorem}
			\begin{figure}[h!]
				\centering
				\includegraphics[width=0.4\textwidth]{figures/fig_laurents-theorem}
			\end{figure}
			Let $C_1$ and $C_2$ be concentric circles of radii $R_1$ and $R_2$, respectively, and center at $a$. Suppose that $f(z)$ is single-valued and analytic on $C_1$ and $C_2$ and, in the ring-shaped region $\mathcal{R}$ between $C_1$ and $C_2$. Let $z=a+h$ be any point in $\mathcal{R}$. Then we have
			\begin{align*}
				f(z)&=\cdots+\dfrac{a_{-3}}{(z-a)^3}+\dfrac{a_{-2}}{(z-a)^2}+\dfrac{a_{-1}}{z-a}+a_0+a_1(z-a)+a_2(z-a)^2+a_3(z-a)^3+\cdots
			\end{align*}
			where
			\begin{align*}
				a_n&=\dfrac{1}{2\pi i}\oint_{C_1}{\dfrac{f(z)\dd{z}}{(z-a)^{n+1}}};\;n=0,1,2,\dots \\
				a_{-n}&=\dfrac{1}{2\pi i}\oint_{C_1}{(z-a)^{n-1}f(z)\dd{z}};\;n=1,2,3,\dots
			\end{align*}
			with $C_1$ and $C_2$ being traversed in the positive direction with respect to their interiors.
	
	%\subsubsection{Residue theorem}